% !TEX program = xelatex
% Example: Subfigures with subcaption
\documentclass[11pt, a4paper]{article}
\usepackage{graphicx}
\usepackage{subcaption}

\begin{document}

\section*{الأشكالُ الفرعية}

% شكلان جنباً إلى جنب
\begin{figure}[htbp]
  \centering
  \begin{subfigure}{0.45\textwidth}
    \centering
    \includegraphics[width=\linewidth]{example-image-a}
    \caption{النتائجُ التجريبيةُ}
    \label{fig:exp}
  \end{subfigure}
  \hfill
  \begin{subfigure}{0.45\textwidth}
    \centering
    \includegraphics[width=\linewidth]{example-image-b}
    \caption{نتائجُ المحاكاة}
    \label{fig:sim}
  \end{subfigure}
  \caption{مقارنةٌ بينَ النتائجِ التجريبيةِ والمحاكاة}
  \label{fig:comparison}
\end{figure}

% إحالة لشكلٍ فرعيٍّ محدّد
الشكلُ \ref{fig:exp} يُظهرُ النتائجَ التجريبيةَ، بينما
الشكلُ \ref{fig:sim} يُظهرُ نتائجَ المحاكاة.

\end{document}
